%% notes-tex/genesis/genesis-proposal-redacted/sections/milestones.tex
%% [[genesis-proposal-redacted]]
%% https://github.com/Mjvolk3/torchcell/tree/main/notes-tex/genesis/genesis-proposal-redacted/sections/milestones.tex
%% SOURCE: hand-transcribed from the proposal docx, "Milestones in the Nine Months"
%% and its Table 1. The table is a hand-authored verbatim copy of the submission's
%% table; nothing in it is redacted.

\section{Milestones in the nine months\secstatus{todo}}\label{sec:milestones}

This proposal is structured to establish milestones towards building MEFMs with a
6-month decision point aligned with the go/no-go metrics
(Table~\ref{tab:gonogo}). These goals collectively organize the research program
into stages that help address gaps in the development of MEFMs. These include
scaling data across organisms and genotypes, designing a metabolic representation
module, and demonstrating the real-world use of these models by designing strains.
Following the development of the first MEFMs, these models will be used in
metabolic engineering applications to generate data from recommended strains, which
will be fed back into the models to complete the first continuous improvement loop
for MEFMs.

In the first 6 months, there will be a computational track (i--iii, v) adding new
data to the current TorchCell database, developing TorchCell CGT models for
\org{E. coli}, and grounding virtual cell models to respect mass balance and
thermodynamics; and the experimental track (iv) where a protocol will be developed
for an RNA-barcoded YKOC experiment in \org{I. orientalis} where the subsequent
atlas will be published. Improved TorchCell CGT models for \org{S. cerevisiae} and
\org{E. coli} will be finetuned (v) using published data for \org{P. putida} and
preliminary in-house data for \org{I. orientalis}. These models will be used for
strain design recommendations for 3-HP production in \org{I. orientalis} and
isoprenol production in \org{P. putida}. At the 6-month decision point, success
will be evaluated using quantitative thresholds aligned with the decision matrix
(Table~\ref{tab:gonogo}). Achieving these criteria (i--v) will demonstrate the
scaling of TorchCell's database and experimental ontology across diverse organisms
and phenotypes, and that physically grounded MEFMs can be used for strain design
across eukaryotes and prokaryotes.

\begin{table}[H]
  \centering\footnotesize
  \caption{Six-month go/no-go decision metrics and AI advantage criteria, as
  submitted (the proposal's Table 1).}
  \label{tab:gonogo}
  \begin{tabular}{L{30mm}L{62mm}L{62mm}L{12mm}}
    \toprule
    \textbf{Metric} & \textbf{Go} & \textbf{AI advantage} & \textbf{Months} \\
    \midrule
    i) TorchCell Database Expansion &
    10 bacteria and 10 yeasts plus key organisms \org{E. coli}, \org{P. putida},
    and \org{I. orientalis} added to the TorchCell database with $\geq$ 1,000
    experiment instances and $\geq$ 3 phenotypes each. &
    Shared experimental ontology with TorchCell dataloaders converts heterogeneous
    data $\geq$ 20 organisms into benchmark-ready DL datasets for building MEFMs. &
    1--6 \\
    \midrule
    ii) TorchCell CGT \org{E. coli} &
    TorchCell CGT benchmarked on \org{E. coli} demonstrating phenotype prediction
    exceeding FBA and linear baselines $\geq$ 20\%. &
    $\geq$ 20\% improvement over FBA and linear baselines demonstrates multimodal
    representations capture biological nonlinearities beyond shallower methods. &
    2--4 \\
    \midrule
    iii) Metabolic Grounding &
    Addition of metabolism module to TorchCell CGT so that $\geq$ 95\% of top 100
    strains don't violate mass balance or $\Delta G$ feasibility. &
    Physical constraints filter implausible strain designs and are independent of
    training-data scale, critical for non-model organisms. &
    1--6 \\
    \midrule
    iv) \org{I. orientalis} Atlas &
    Developed protocol for \org{I. orientalis} RNA-barcoded YKOC across $\geq$ 3
    environmental conditions yielding $\geq$ 5,000 perturbation-condition pairs. &
    First genome-$\times$-environment perturbation atlas in \org{I. orientalis}
    establishing a template for scaling across organisms and providing highly dense
    training signal for this organism of interest. &
    1--9 \\
    \midrule
    v) Finetune TorchCell CGT &
    TorchCell CGT finetuned from \org{S. cerevisiae} to \org{I. orientalis}, and
    from \org{E. coli} to \org{P. putida} matches from scratch baseline Pearson
    with $\geq$ 5$\times$ fewer target labels. &
    5$\times$ sample efficiency on cross-organism transfer quantifies the returns
    of adding new host organism data and provides scaling evidence for expanding
    MEFMs. &
    4--6 \\
    \midrule
    vi) TorchCell CGT Foundation Model Training Expansion &
    TorchCell CGT models trained across $\geq$ 10 yeasts and 10 $\geq$ bacteria
    with leave-one-organism-out Pearson $\geq$ 0.3 and reported transfer exponents
    with 95\% CIs. &
    Scaling exponents with uncertainty quantify the return per added organism and
    phenotype, providing measured scaling behavior for MEFM training. &
    6--9 \\
    \midrule
    vii) TorchCell CGT Metabolic Engineering Applications &
    $\geq$ 1 strain showing ($p < 0.05$, $n \geq 3$) titer improvement over base
    strain for 3-HP production in \org{I. orientalis} and isoprenol production in
    \org{P. putida}, verified by high-throughput mass spec. &
    Automated biofoundries scale experimental DBTL throughput while feeding results
    back to MEFM, demonstrating computational and experimental scaling for future
    strain engineering campaigns. &
    6--9 \\
    \bottomrule
  \end{tabular}
\end{table}

During months 6--9 TorchCell models will be retrained (vi) on data within the same
domain of life to lay the blueprint for building MEFMs for broad phenotype
prediction, but specifically industrial metabolic engineering application. Results
from MEFMs trained across domains of life will be used to refine strain designs,
while the first designs from finetuning are being validated (vii).
